\section{Progress-driven discovery before planar refinement}
\label{sec:planar-adaptive}

Complete planar enumeration and first-disc discovery have different cost
profiles. The preceding release prepares every minimum cell and crossing
before yielding a ray. Its small two-cap discovery cases are slower than
the former arrangement iterator, even though complete enumeration is much
faster. The new discovery policy responds to observed progress: test the
original projective Q-corners first, and build the complete minimum overlay
only when that prelude has not found a disc. Its complete fallback and
proof formats are preserved.

\subsection{A cheap prefix of the complete ray set}

Let $P$ be the normalized nonnegative quadrilateral section. Every corner
of $P$ is extreme in any minimum cell containing it. Canonical lifting is
a linear isomorphism from that cell's cone to a standard coordinate face,
so its corner lift is a standard extreme ray. This is the face theorem of
Sections~\ref{sec:minimum-envelopes-native} and
\ref{sec:planar-overlay-refinements}; it does not depend on the Euler
convexity hypothesis about vertex links. A Q-corner is therefore a safe
first candidate even when the source is not an authenticated knot exterior.
It still needs the existing exact essential-disc certificate before any
positive result is returned.

For matching nullity three, construct the source Q-chart and its polygon,
segment, point or empty section without constructing minimum cells.
Project the retained triangle forms once. There are $O(k)$ forms, so the
preparation costs $O(k^2)$ rational operations after kernel construction.
Order domain corners by decreasing first allowed quadrilateral coordinate,
with exact chart-coordinate ties. This is a deterministic candidate policy,
not a theorem predicting essentiality. The first-corner preference is
independent of the Fibonacci fixture parameters. Sorting uses constant-size
keys in $O(k\log(k+2))$ comparisons. One full canonical lift costs
$O(t+k)$ arithmetic operations; the first candidate is thus available in
$O(1+t+k^2)$ post-kernel operations, with rational-bit, component-query,
independent-replay and output costs separately charged.

\begin{proposition}[Complete adaptive discovery stream]
Yielding all domain-corner lifts, then all remaining vertices of the common
minimum subdivision, gives each non-link standard extreme ray exactly once.
If a consumer stops at a successful corner, no minimum-overlay construction
is required. If it resumes after the prelude, the complete stream has the
same ray set as full planar enumeration.
\end{proposition}
\begin{proof}
Domain corners are vertices of every subdivision and their lifts are
standard-extreme by the coordinate-face correspondence. Build the common
minimum subdivision using the already prepared chart and projected forms.
Its vertices are exactly the complete standard-ray set. Remove only the
exact domain points already yielded; normalization meets each nonzero ray
once, so neither duplicate rays nor omissions result. Empty and singleton
sections complete at the prelude. A segment resumes through its exact
minimum breaks; a polygon resumes through all cells and cross-group
intersections. The generator's control flow builds this refinement only
after the consumer asks for a candidate beyond the prefix.
\end{proof}

The full geometric arithmetic bound stays
$O(1+k^3+(t+k)R)$ under the deterministic realization described earlier.
Every Q-corner is among the $R$ output rays and is lifted once. The prefix
shares its chart and projected forms with the tail. It does not pay a
second full preparation. This is a worst-case preservation result and an
early-exit opportunity, not a uniform speedup theorem. Python hash-set
collision costs retain the qualifications of the preceding section.

\subsection{Complete enumeration at logarithmic nullity}

The general arrangement fallback also gives a useful local complexity
statement beyond the planar gate. This is elementary accounting for the
existing implementation, without a novelty claim.

\begin{proposition}[Complete low-nullity sector enumeration]
For a supplied compatible sector with $k$ allowed types, $t$ tetrahedra,
matching nullity $d$, and encoded source size $B$, the general standard-ray
arrangement algorithm has bit cost
\[
 \operatorname{poly}(t+k+B)(1+k)^{O(d)}.
\]
The higher-dimensional automatic support/arrangement selector preserves
this bound. Hence polynomial source size and logarithmic matching nullity
give quasi-polynomial complete enumeration within one supplied sector.
\end{proposition}
\begin{proof}
The contracted matching model has at most $8k$ retained triangle classes.
If $p_v$ classes belong to vertex group $v$, the candidate hyperplanes
comprise $k$ coordinate forms and at most $\sum_v\binom{p_v}{2}=O(k^2)$
potential differences, projected to the $d$-dimensional matching kernel.
Identically zero and duplicate forms can be dropped. An extreme ray of a
minimum-cell cone is determined by $d-1$ independent active hyperplanes
from this list. This remains true when the feasible cone has smaller
dimension: its forced equalities are among the active constraints.
Thus testing all $\binom{O(k^2)}{d-1}$ systems, followed by the exact
standard-extremality filter, is complete. Nullity zero and empty sectors
have their explicit no-ray branches.

Each candidate requires polynomial-size exact linear algebra, canonical
lifting, source-coordinate validation and output. The bounded integral
matching coefficients and determinant bounds control elimination bits;
the inherited $4^s$ and $4^{s-1}$ coordinate bounds for actual support $s$
give $O(k)$ bits per primitive output coordinate. No elementary pieces
are expanded. The automatic support alternative has at most its computed
mixed-support count in candidate bases and is selected only when that
count is below the arrangement count. Its per-candidate elimination is
also polynomial. Computing the two combinatorial bounds uses polynomially
many integer operations on polynomial-bit integers. Multiplying the
candidate count by these polynomial costs gives the result.
Even quadratic worst-case duplicate-table work only squares the candidate
count and stays within $(1+k)^{O(d)}$; no constant-time hashing assumption
is needed for this broader complexity class.
\end{proof}

This strengthens the local enumeration accounting beyond the Euler-only
logarithmic-nullity observation in report 81. It supplies neither a
quasi-polynomial complete sector producer nor a bound for source hierarchy
resets and essential-disc testing across all knot inputs. Those additional
requirements are still needed for a general recognition theorem.

For the two-cap Fibonacci family, the first-corner behavior is uniform,
not only an observation on the timed sizes. Its Q-cone is
$(x,y,z)\in\mathbb R_{\ge0}^3$. Only the $x$-corner has nonzero first
allowed coordinate $q_0=F_nx$; the two cap corners have $q_0=0$.
The decreasing-coordinate policy therefore selects $x=1,y=z=0$ first
after primitive normalization, since the last base coordinate is $F_1=1$.
By the all-size component theorem of Section~\ref{sec:planar-sectors-native},
this is exactly one essential meridian disc. A complete exact disc
consumer thus succeeds on its first candidate for every $n\ge1$ without
requesting a minimum overlay. Cooperative interruption or an observer
resource limit can still produce an inconclusive query. This family theorem
does not predict first-corner success for arbitrary triangulations.

\subsection{One allowance, ordinary positive and negative certificates}

The native generator is \code{sector\_planar\_discovery\_rays}. It uses
one actual-ray allowance across the prefix and tail: its counter is not
reset when the implementation switches to refinement. The already tested
corner points are excluded before any further coordinate lift. A cap
produces \code{INCONCLUSIVE} without an exhaustion certificate. The
cooperative callback covers section construction, lifting, resumed
refinement and the existing component checks; callable objects with false
Boolean value still receive cancellation calls.

Automatic STANDARD discovery selects this policy at matching nullity
three. The other automatic dimensions and explicit methods retain their
previous behavior. Complete enumeration and planar coverage APIs retain
their previous canonical order and statistics; they have no reason to
prioritize a potential first disc.

The first candidate is inspected by the maintained compressed disc
observer and independently replayed. A success uses the unchanged
\code{normal-sector-witness-v1} source-bound proof. Exhaustion uses the
unchanged \code{normal-sector-exhaustion-v1} with phase STANDARD and every
ray's Euler value and required negative disc proof. Its independent dense
checker compares a set, so a different traversal order changes no coverage
claim. In particular, failing the corner prelude does not become a negative
disc result. We deliberately do not install the proposed aggregated
Euler-only negative precheck here: this stream keeps full standard-ray
coverage and its established portable certificate contract.

\subsection{Reusing a kernel after a prior Q phase}

The bounded-support outer search already tries a complete Q phase before
requesting STANDARD discovery on a \code{POSITIVE\_EULER\_ONLY} result.
Starting a new corner prelude there would repeat work. The search now builds
one producer matching kernel per selected sector and passes it through
both phases. Its internal STANDARD continuation skips the new prelude
because the Q phase has already tested those candidates; it uses the
ordinary complete standard iterator. Per-phase candidate allowances keep
their previous meaning. The established STANDARD checks still include
their own corner entries; this change removes the redundant producer build
and avoids adding another discovery prelude. The source/kernel is read-only throughout this
query. Independent proof checkers continue to reconstruct geometry from
the actual supplied source; cached producer geometry is not substituted
for source authentication or verification.

This removes a redundant producer build for each sector requiring a
second phase. It does not yet share preparation across different sectors
and does not change which supports are enumerated. The outer negative
status remains restricted to the declared support allowance, with no
unqualified knot verdict. This supplied-triangulation research API is
separate from the default diagram-recognition portfolio.

\subsection{Validation and measurements}
All 1,396 maintained tests pass in 227.198 seconds. The focused 48-test
run passes in 1.719 seconds. Tests forbid overlay construction on an early
disc, compare complete unique adaptive streams on every support of the
three-tetrahedron two-cap source, replay a seven-ray negative result,
remove each ray from its certificate, stop between stages, and check one
shared cap at limits zero, one, three, six and seven. A genuine capped
trefoil control verifies one producer kernel per selected support and
identity reuse across failed Q and STANDARD phases.

The new native audit exactly reproduces the incumbent's complete canonical
ray sequences and geometry statistics on all 9,807 selected eligible
sectors. Adaptive complete streams give the same ray sets, without
duplicates, on all those queries. On every one of the 463 selected
nullity-three sectors, old and new discovery statuses match and the new
certificate independently replays. All 44 observed positive queries finish
in the corner prelude. All 419 negative proofs replay, with 415 requiring
the full overlay and four ending on degenerate sections. Their transcript
entry sets exactly equal the incumbent's. Eight fresh Regina full
enumerations again match the frozen reference. This is the previously
declared selection across 48 triangulations, not all supports on all
possible inputs. All 589 captured source pins remain fixed during audit
and timing.

The baseline is \code{89d3e8bee}: both its native sector module and planar
module are frozen, with only the sector module's lazy import redirected to
the frozen planar module. Other imported dependencies are checked unchanged
at the runtime release. This prevents the incumbent timing from silently
using the new planner. Five serial interleaved four-arm rounds complete
280 measured calls and 56 warmups. Source fixtures are prepared before
timing. Discovery includes full kernel construction, actual disc or
exhaustion proof, independent replay and complete serialization. Individual
old/new witnesses can differ, so both are checked and their byte lengths
are retained. No partial or capped run contributes a ratio.

\begin{samepage}
\noindent\textbf{Complete supplied-sector discovery and proof replay.}
\begin{center}
\small
\begin{tabular}{lrrrrr}
\toprule
Input & old ms & new ms & old/new & old A/A & new A/A\\
\midrule
Base 1 & 6.683 & 2.498 & 2.684 & 1.013 & 1.011\\
Base 4 & 12.500 & 4.966 & 2.521 & 1.009 & 0.976\\
Base 8 & 20.785 & 8.848 & 2.363 & 1.012 & 0.991\\
Base 16 & 42.106 & 23.028 & 1.874 & 0.996 & 1.023\\
Base 32 & 107.636 & 63.952 & 1.641 & 1.014 & 1.005\\
Complete negative & 28.560 & 28.914 & 0.994 & 0.995 & 0.997\\
Empty sector & 1.410 & 1.412 & 1.000 & 0.999 & 0.991\\
\bottomrule
\end{tabular}
\end{center}

\end{samepage}

The size-one query drops from about 6.683 to 2.498 milliseconds, a paired
2.684-fold gain. At size 32 it drops from about 107.636 to 63.952
milliseconds, a paired 1.641-fold gain. The full negative and empty-sector
controls stay near parity. These ratios compare directly with the preceding
planar release on the same host; they are not assembled from timings in
different releases. The first-corner family proof above explains the
successful prefix, while the complete negative control exercises the
resumed fallback. There is no claim of a uniform first-corner success rate.

\begin{samepage}
\noindent\textbf{Unchanged lower-nullity and complete-enumeration controls.}
\begin{center}
\small
\begin{tabular}{lrrrrr}
\toprule
Input & old ms & new ms & old/new & old A/A & new A/A\\
\midrule
One cap 4 & 3.440 & 3.452 & 0.996 & 0.999 & 1.002\\
One cap 8 & 6.634 & 6.611 & 1.003 & 1.005 & 1.008\\
Enumeration 1 & 3.175 & 3.189 & 0.996 & 0.999 & 1.012\\
Enumeration 8 & 11.268 & 11.180 & 1.004 & 1.013 & 1.004\\
Enumeration 16 & 25.141 & 24.671 & 1.018 & 0.986 & 0.973\\
\bottomrule
\end{tabular}
\end{center}

\end{samepage}

Complete enumeration includes every output coordinate, serialization and
ray digest; portable coverage replay is outside that iterator timing.
It retains its previous output and costs, with paired ratios near one.
Lower-nullity discovery also stays near parity. A/A medians over the whole
run range from 0.973 to 1.023.

\begin{samepage}
\noindent\textbf{Complete bounded-support scans on supplied capped exteriors.}
\begin{center}
\small
\begin{tabular}{lrrrrr}
\toprule
Input & old ms & new ms & old/new & old A/A & new A/A\\
\midrule
Capped trefoil & 33.488 & 31.639 & 1.052 & 0.997 & 0.973\\
Capped figure eight & 64.924 & 63.044 & 1.031 & 1.004 & 1.015\\
\bottomrule
\end{tabular}
\end{center}

\end{samepage}

These outer API calls exhaust every support of size at most one. Native
candidate checks are included; this API does not emit a portable proof
of the entire support-family exhaustion, and such replay is not claimed
inside these times. Both arms return identical restricted-scope records.
The first paired gains are about 1.052 and 1.031. Because these effects
are small, a separate fifteen-round four-arm repeat completes 120 measured
calls and eight warmups. Its gains are 1.052 and 1.041, with A/A medians
between 0.996 and 1.005. Both runs are retained. The structural result is
one producer build per support; the small timing estimates do not imply
the same gains on every triangulation or in the diagram portfolio.

\begin{sloppypar}
The maintained driver is \code{normal\_orbit\_research.planar\_adaptive},
with \code{audit --fresh-regina} and \code{benchmark --rounds 5} modes.
Raw evidence and the longer repeat are in \code{data/planar-adaptive-*}.
Publication replay checks frozen source versions and saved certificates
with the producer routines disabled. It retains the existing thirteen
coverage proofs, all 463 discovery certificates and the 44 incumbent
positive witnesses. General recognition and complete sector-family
production remain separate unproved requirements.
\end{sloppypar}

